\documentclass{article}
\usepackage[portuguese]{babel}
\usepackage{macros}
\title{Lista 4 - MAT01050}
\author{Leonardo Suarez}
\date{\today}

\begin{document}
\maketitle

\begin{unprblm}
    Seja \(A = \begin{bmatrix}
        1 & 1 \\
        2 & 3 \\
        0 & 1
    \end{bmatrix}\) e \(b = \begin{bmatrix}
        2 \\
        5 \\
        1
    \end{bmatrix}\).
    \begin{itemize}
        \item Encontre a solução única de MQ usando (1) \(x = A^\dagger b\), (2) as equações normais, (3) os métodos de fatoração QR com Householder e Givens, e (4) os métodos CGS e MGS.
        \item Determine \(\kappa(A)\).
        \item Mostre que para esse problema, a sensitividade do problema de MQ, quando apenas \(A\) é perturbada, depende de \(\kappa(A)\).
        \item Seja \(E = 10^{-4}A\) e seja \(\overline{x} = \overline{A}^\dagger b\) , onde \(\overline{A} = A + E\). Calcule \(\frac{\Vert \overline{x}-x \Vert}{\Vert x \Vert}\) e verifique a desigualdade \(\frac{\Vert \overline{x}-x \Vert}{\Vert x \Vert} \leq \kappa(A) \frac{\Vert E \Vert}{\Vert A \Vert}\).
        \item Encontre \(r\) e \(\hat{r}\) e verifique a desigualdade \(\frac{\Vert \hat{r} - r \Vert}{\Vert r \Vert} \leq \kappa(A) \frac{\Vert E \Vert}{\Vert A \Vert}\).
    \end{itemize}
\end{unprblm}

\begin{unprblm}

\end{unprblm}

\begin{unprblm}

\end{unprblm}

\end{document}